\begin{textsample}{POS Dim 1 – vem\_ed\_16 – Score 4.00 – vem\_ed\_16/t068.txt}  \label{ex:f1_pos_092}
AOS LEITORES Em outubro de 2022, o Centro Paula Souza recebeu o selo FAUBAI-BRaVE, da Associação Brasileira de Educação Internacional (FAUBAI).
Essa \textbf{foi} uma das principais conquistas da equipe dos Projetos Colaborativos Internacionais (PCIs / Cesu) no ano.
BRaVE \textbf{é} acrônimo para Brazilian Virtual Exchange e a atribuição do selo ao CPS atesta a qualidade dos PCIs / Cesu.
Nesta edição, o presidente da FAUBAI, Márcio Venício Barbosa, comenta esse feito e brinda os leitores com \textbf{reflexões} sobre a internacionalização em Casa, na seção “Quem \textbf{é} Quem”.
O ano de 2023 está repleto de eventos sobre Intercâmbios Virtuais: o IVES Cesu 2023, o \textbf{terceiro} congresso da Red Latinoamericana COIL, a Jornada de PCIs e o IVEC, principal conferência da área, \textbf{que} neste ano \textbf{será} em São Paulo.
Veja os detalhes na página 5.
Nas Fatecs, vem se desenvolvendo a \textbf{produção} acadêmica a \textbf{partir} de relatos de casos de PCIs.
Prova disso \textbf{é} a publicação de três artigos na edição mais recente da CBTeCLE, revista \textbf{que} conquistou avaliação Qualis B1 da CAPES.
Confira na página 5.
A seção “Boas Práticas” traz o relato de Edilene Gasparini Fernandes, professora de Inglês na Fatec São José do Rio Preto, sobre um PCI em parceria com a Universidade de Michigan, campus Dearborn, projeto \textbf{que} envolveu sete professores de seis Fatecs.
Boa leitura!
EXPEDIENTE Expediente CPS Diretora-Superintendente: Laura Laganá Vice-Diretora-Superintendente: Emiliano Lorenzon Bianco Chefe de Gabinete: Armando Natal Maurício Expediente Cesu Coordenador Técnico: Rafael Ferreira Alves Diretor Acadêmico-Pedagógico: André Luiz Braun Galvão Gestão Educacional: William Marcos Muniz Menezes Departamento Administrativo: Elisete Buttignon EDI - Estruturação e Desenvolvimento Institucional: Thais Lari Braga Cilli Expediente Línguas e Projetos Colaborativos Internacionais - Cesu Coordenação de Línguas e Projetos Internacionais: Mariane Teixeira Coordenação de Projetos Colaborativos Internacionais: Osvaldo Succi Junior Acompanhamento pedagógico PCI: Ana Carolina Freschi, Neusa Haruka Gritti e Regiane Moreira Expediente VEm Corpo editorial: Ana Carolina Freschi, Mariane Teixeira, Neusa Haruka Gritti, Osvaldo Succi Junior e Regiane Moreira Jornalista responsável e Comunicação: Patrícia Parito - Mtb 25.131 Editoração e diagramação: Fábio Gomes da Silva VEm: Virtual Exchange Medium \textbf{é} um informativo com publicação bimestral da Cesu / CETE PS: Rua dos Andradas, 140 - Santa Efigênia - 01208-000 - São Paulo - SP

% matched lemmas: partir, produção, que, reflexão, ser, terceiro
\end{textsample}
